\section{Equilibrium results}\label{sec:results}

\subsection{Competition creates competition}\label{sec:ccc}

The main result compares three economies that differ only in incumbent strength. In the first, the price carries
no information and only a zero-cost challenger prepares. In the second, the price is informative and a purchase
recruits a high-cost challenger. In the third, the price stays informative and entry returns to its floor. Write
\(F_0\), \(F_1\) and \(F_2\) for the three incumbent distributions, \(B_i=B_{F_i}\), \(\Delta_i=\DT(F_i)\),
\(M=M(\pi)\), \(m=m(\pi)\) and \(\kappa=\kappa(\pi)\).

\begin{proposition}[Competition creates competition; analytical]\label{prop:ccc}
Let the classes satisfy Assumption~\ref{ass:classes}. Let \(F_1\) be a continuous first-order stochastic
strengthening of the continuous incumbent distribution \(F_0\), and let \(F_2\) be one of \(F_1\). Fix
\(0<\pi<1\), \(0<\rho<1\), \(k>0\) and a slot \((\lambda,\beta)\). Suppose
\begin{align*}
&\text{(H1)}\qquad B_0(\pi)<c_H<B_1(M),\\
&\text{(H2)}\qquad \Delta_0<k<\rho\,\kappa\,\Delta_1,\\
&\text{(H3)}\qquad B_2(M)<c_H.
\end{align*}
Then each economy has a unique equilibrium outcome.
\begin{enumerate}[label=(\roman*),leftmargin=2em]
\item \emph{Weak incumbent \(F_0\).} The investor does not trade: \(\sigma_H=\sigma_L=\delta_0\). Every posterior
equals \(\pi\), entry equals \(\rho\) after every order, and every order has the price
\(t_0+\rho\,[t_L-t_0+\Delta_0\pi]\). The price carries no information.
\item \emph{Strong incumbent \(F_1\).} The investor trades in both directions: \(\sigma_H=\delta_1\) and
\(\sigma_L=\delta_{-1}\). The posteriors after the orders \((-1,0,1)\) are \((m,\pi,M)\). Only a purchase recruits
the high-cost challenger, so entry after the three orders is \((\rho,\rho,1)\), and \(P(-1)<P(0)<P(1)\). Entry is
\[
\Ent_1=\rho+(1-\rho)(b+\pi\lambda)>\rho.
\]
Entry given class \(H\) is \(e_H=\rho+(1-\rho)(b+\lambda)\). High-class ownership rises from \(\pi\rho W_H(F_0)\)
to \(\pi e_HW_H(F_1)\) if and only if
\[
[\rho+(1-\rho)(b+\lambda)]\,W_H(F_1)>\rho\,W_H(F_0).
\]
\item \emph{Very strong incumbent \(F_2\).} The investor again trades in both directions, with posteriors
\((m,\pi,M)\). Entry equals \(\rho\) after every order. The three prices differ, so the price stays informative
while entry returns to \(\rho\).
\end{enumerate}
\end{proposition}

Uniqueness covers every mixed order profile on the menu \(\{-1,0,1\}\) and every unilateral deviation. Every
investor comparison in the proof is strict, and every equilibrium posterior lies strictly on one side of the
threshold. No tie rule is therefore used, and the same unique outcome holds under all four pairs of tie rules.
Appendix~\ref{app:prop2} states the three equilibrium tuples and checks each condition of
Definition~\ref{def:equilibrium} against them.

The hypotheses fix trading before any order profile is known. By Lemma~\ref{lem:trading}, every correct order
earns at most \(\Delta_0-k<0\) in the weak economy and at least \(\rho\kappa\Delta_1-k>0\) in the strong one. The two
sides of (H1) set entry. Gross profit at the prior falls short of the high cost at \(F_0\), and gross profit at the
top posterior exceeds it at \(F_1\). Hypothesis (H3) stops recruitment at \(F_2\) (Corollary~\ref{cor:exclusion}).
Appendix~\ref{app:prop2} shows that the proof uses each hypothesis once and that (H1) cannot be weakened to
\(B_1(\pi)<c_H\).

The trading window in (H2) is nonempty exactly when \(\Delta_0/\Delta_1<\rho\kappa\). This condition involves the
two incumbent distributions, the classes, \(\rho\), the slot and \(\pi\), and no unit of money. A prior away from
one half narrows the window, since \(\kappa(\pi)\) falls. For uniform incumbents, (H1)--(H3) are strict
inequalities between continuous functions of \((r_0,r_1,r_2,c_H,k,\rho,\pi,\lambda,\beta)\). They define an open
set. Section~\ref{sec:benchmark} gives its slice through the benchmark.

The additional entry favors the high class: class \(H\) places a purchase with probability \(b+\lambda\), class
\(L\) with probability \(b\). Ownership needs its own condition when the high class can lose to the incumbent,
because a stronger incumbent then lowers \(W_H\). The condition in part~(ii) holds whenever
\((1+\rho)W_H(F_1)\ge2\rho W_H(F_0)>0\) and \(b+\lambda>1/2\). In the point-mass benchmark \(W_H=1\), so ownership
rises.

Part~(iii) completes a rise and a later fall of entry with incumbent strength. Against the very strong incumbent,
even the top posterior fails to justify high-cost preparation, while trading stays informative.
Section~\ref{sec:coexistence} describes the strengths between the three economies.

\subsection{The benchmark and the fixed-information control}\label{sec:benchmark}

The benchmark, which I call \(B\), has point-mass classes \(h=\reg{prim_h}\) and \(\ell=\reg{prim_ell}\) and the
prior \(\pi=\reg{prim_pi}\). The reserve is \(p=\reg{prim_p}\), the zero-cost probability is \(\rho=\reg{prim_rho}\),
the high cost is \(c_H=\reg{prim_cH}\), and the trading cost is \(k=\reg{prim_k}\). The incumbent is uniform on
\([0,r]\), at the strengths \(r_0=\reg{prim_r0}\), \(r_1=\reg{prim_r1}\) and \(r_2=\reg{prim_r2}\). Money is
measured in units of the low value.

The slot has two primitives. The largest likelihood ratio of an order is \(1+u=\eu\), where \(\eu\) is Euler's
number; I print it upright to separate it from entry \(e(x)\). The liquidity probability is \(\beta=\reg{beta_B}\), an exact decimal.
The presence probability follows: \(\lambda=u\beta/(1+u\beta)=\reg{lambda_B}\), and \(b=\reg{b_B}\). Every
computation keeps \(\lambda\) unrounded, and every statement that depends on \(\eu\) uses interval arithmetic. The
posterior bounds at the prior are \(M(1/2)=\eu/(1+\eu)=\reg{M_half_B}\) and \(m(1/2)=\reg{m_half_B}\).
Appendix~\ref{app:parameters} lists these declarations.

\begin{table}[tbp]
\centering\small\singlespacing
\caption{Acquisition payoffs in the benchmark.}\label{tab:payoffs}
\begin{tabular}{@{}lccc@{}}
\toprule
 & \(r_0=\reg{prim_r0}\) & \(r_1=\reg{prim_r1}\) & \(r_2=\reg{prim_r2}\)\\
\midrule
Proceeds without entry, \(t_0\) & \reg{p2_base_r0_t0} & \reg{p2_base_r1_t0} & \reg{p2_base_r2_t0}\\
Proceeds with a high-class challenger, \(t_H\) & \reg{p2_base_r0_th} & \reg{p2_base_r1_th} & \reg{p2_base_r2_th}\\
Proceeds with a low-class challenger, \(t_L\) & \reg{p2_base_r0_tl} & \reg{p2_base_r1_tl} & \reg{p2_base_r2_tl}\\
Target-payoff spread, \(\DT\) & \reg{Delta_r0} & \reg{Delta_r1} & \reg{Delta_r2}\\
Gross profit of class \(H\), \(g_H\) & \reg{p2_base_r0_gh} & \reg{p2_base_r1_gh} & \reg{p2_base_r2_gh}\\
Gross profit of class \(L\), \(g_L\) & \reg{p2_base_r0_gl} & \reg{p2_base_r1_gl} & \reg{p2_base_r2_gl}\\
Gross profit at the prior, \(B_r(1/2)\) & \reg{Bprior_r0} & \reg{Bprior_r1_B} & \reg{Bprior_r2}\\
Preparation threshold, \(\tau(r)\) & \reg{tau_r0} & \reg{tau_r1} & \reg{tau_r2}\\
\bottomrule
\end{tabular}
\par\medskip
\begin{minipage}{0.92\linewidth}\footnotesize
\textit{Notes.} Point-mass classes \(h\) and \(\ell\), reserve \(p\), incumbent uniform on \([0,r]\). The entries
follow from equation~\eqref{eq:payoffs}. They do not depend on the slot.
\end{minipage}
\end{table}

Table~\ref{tab:payoffs} reports the acquisition payoffs. From \(r_0\) to \(r_1\), gross acquisition profit at the
prior falls from \(\reg{Bprior_r0}\) to \(\reg{Bprior_r1_B}\), while the spread rises from \(\reg{Delta_r0}\) to
\(\reg{Delta_r1}\). The benchmark satisfies (H1)--(H3), so Proposition~\ref{prop:ccc} applies. Entry rises from
\(\reg{entry_r0}\) to \(\reg{entry_r1_B}\), and high-class ownership rises from \(\reg{own_r0}\) to
\(\reg{own_r1_B}\). At \(r_2\), entry returns to \(\reg{entry_r2}\), while the three prices still differ.
Table~\ref{tab:outcomes} lists the three equilibrium tuples.

The admissible region is a result. Proposition~\ref{prop:region} in Appendix~\ref{app:prop2} gives it exactly at the
benchmark auction. At the prior \(1/2\), (H1)--(H3) hold if and only if \(\reg{u_lo_frac}<u<\reg{u_hi_frac}\).
Equivalently, the top posterior lies between the thresholds \(\tau(r_1)=\reg{M_lo_frac}\) and
\(\tau(r_2)=\reg{M_hi_frac}\). Outside this band a conclusion of Proposition~\ref{prop:ccc} fails. With \(1+u=\eu\),
the hypotheses hold exactly for priors between \(\approx\reg{prior_lo_B}\) and \(\approx\reg{prior_hi_B}\).
Table~\ref{tab:region} gives the region in decimals, with the benchmark point and its distance to each boundary.

\begin{table}[tbp]
\centering\small\singlespacing
\caption{The admissible region at the benchmark auction.}\label{tab:region}
\begin{tabular}{@{}lcccc@{}}
\toprule
Parameter & Region & Benchmark & \multicolumn{2}{c}{Margin to the boundary}\\
\cmidrule(l){4-5}
 & & & lower & upper\\
\midrule
Mass ratio \(u\) & \((\reg{u_lo},\reg{u_hi})\) & \reg{u_B} & \reg{margin_u_lo_B} & \reg{margin_u_hi_B}\\
Top posterior \(M(1/2)\) & \((\reg{M_lo},\reg{M_hi})\) & \reg{M_half_B} & \reg{margin_recruit_post_B} & \reg{margin_return_post_B}\\
Presence \(\lambda\) at \(\beta=\reg{beta_B}\) & \((\reg{lambda_lo_B},\reg{lambda_hi_B})\) & \reg{lambda_B} & \reg{margin_lambda_lo_B} & \reg{margin_lambda_hi_B}\\
Presence \(\lambda\) at \(\beta=1/3\) & \((\reg{lambda_lo_third},\reg{lambda_hi_third})\) & & & \\
Prior \(\pi\), all three parts & \((\reg{prior_lo_B},\reg{prior_hi_B})\) & \reg{prim_pi} & \reg{margin_prior_lo_B} & \reg{margin_prior_hi_B}\\
Prior \(\pi\), parts (i) and (ii) & \((\reg{prior_lo_B},\reg{prior_ii_hi})\) & \reg{prim_pi} & \reg{margin_prior_lo_B} & \\
Liquidity probability \(\beta\) & \((0,1/3]\) & \reg{beta_B} & & \reg{margin_beta_B}\\
\bottomrule
\end{tabular}
\par\medskip
\begin{minipage}{0.92\linewidth}\footnotesize
\textit{Notes.} Benchmark auction and costs, strengths \(r_0,r_1,r_2\). Rows 1 to 4 fix the prior at \(1/2\);
rows 5 and 6 fix \(1+u=\eu\). Each region is the exact open set on which (H1)--(H3) hold, or (H1)--(H2) in row 6.
Decimal endpoints are rounded to six places; Proposition~\ref{prop:region} gives the exact endpoints.
\end{minipage}
\end{table}

The thinnest margins measure one distance. The benchmark top posterior lies \(\reg{margin_return_post_B}\) below
\(\tau(r_2)\). In cost units, \(c_H-B_{r_2}(M)=\reg{margin_return_B}\), and \(r_2\) lies \(\reg{gap_r2_rC_B}\)
above the ceiling \(r_C=\reg{rC_B}\), the strength at which \(\tau(r)=M(1/2)\). Every hypothesis has a positive
margin at the benchmark. Table~\ref{tab:margins} in Appendix~\ref{app:parameters} lists each hypothesis with its
margin.

\begin{figure}[tbp]
\centering
\slotfigure{figures/beliefs.pdf}{6cm}
\caption{Beliefs and the ceiling in the benchmark.}\label{fig:beliefs}
\par\medskip
\begin{minipage}{\linewidth}\footnotesize\singlespacing
\textit{Notes.} Panel (a) plots the threshold belief \(\tau(r)\) against incumbent strength \(r\). The gray band
covers the posteriors that an investor strategy can produce, from \(m(1/2)\) to \(M(1/2)\). The ceiling \(r_C\) is the
strength at which \(\tau(r)\) equals \(M(1/2)\). Panel (b) enlarges the strengths from \(r_1\) to \(r_2\). Its blue band
holds every top posterior \(M(1/2)\) that puts \(r_C\) between \(r_1\) and \(r_2\); the benchmark value lies
\(\reg{margin_return_post_B}\) below the upper edge \(\tau(r_2)\).
\end{minipage}
\end{figure}

Figure~\ref{fig:beliefs} shows the threshold along the path. The prior lies below the threshold at every strength,
so an uninformative price never recruits a high-cost challenger. The threshold rises with strength and meets the top
posterior at \(r_C\), just below \(r_2\). Panel~(b) shows the band of Proposition~\ref{prop:region}(a) for the top
posterior.

Two controls show where the reversal comes from. The first holds the information in the price fixed while the
incumbent strengthens. Under rule C1, entry then does not rise, for every information experiment and every
independent cost law (Proposition~\ref{prop:control} in Appendix~\ref{app:robustness}). A stronger incumbent lowers gross profit at every
belief, so it can only remove a preparation incentive. Entry falls by the probability that the cost lies between gross
profit at the two strengths.

In the slot, the experiment fixes \(\lambda\), \(\beta\), \(\pi\), the cost law and both
order distributions. Lemma~\ref{lem:prices}(c) ensures that the competitive prices reveal the same posterior at each
strength.

\begin{table}[tbp]
\centering\small\singlespacing
\caption{Equilibrium outcomes and information controls in the benchmark.}\label{tab:outcomes}
\begin{tabular}{@{}lccc@{}}
\toprule
 & \(r_0=\reg{prim_r0}\) & \(r_1=\reg{prim_r1}\) & \(r_2=\reg{prim_r2}\)\\
\midrule
\multicolumn{4}{@{}l}{\emph{Panel A. Equilibrium (Proposition~\ref{prop:ccc})}}\\
Orders \((\sigma_H,\sigma_L)\) & \((\delta_0,\delta_0)\) & \((\delta_1,\delta_{-1})\) & \((\delta_1,\delta_{-1})\)\\
Posterior after a sale & \reg{p2_base_r0_mu_minus} & \reg{p2_base_r1_mu_minus} & \reg{p2_base_r2_mu_minus}\\
Posterior after no order & \reg{p2_base_r0_mu_zero} & \reg{p2_base_r1_mu_zero} & \reg{p2_base_r2_mu_zero}\\
Posterior after a purchase & \reg{p2_base_r0_mu_plus} & \reg{p2_base_r1_mu_plus} & \reg{p2_base_r2_mu_plus}\\
High-cost preparation & none & after a purchase & none\\
Price after a sale & \reg{price_r0} & \reg{price_r1_m1_B} & \reg{price_r2_m1_B}\\
Price after no order & \reg{price_r0} & \reg{price_r1_0_B} & \reg{price_r2_0_B}\\
Price after a purchase & \reg{price_r0} & \reg{price_r1_p1_B} & \reg{price_r2_p1_B}\\
Entry \(\Ent\) & \reg{entry_r0} & \reg{entry_r1_B} & \reg{entry_r2}\\
High-class ownership \(\Own_H\) & \reg{own_r0} & \reg{own_r1_B} & \reg{own_r2_B}\\
Expected target proceeds & \reg{p4_base_r0_proceeds} & \reg{p4_base_r1_proceeds} & \reg{p4_base_r2_proceeds}\\
\midrule
\multicolumn{4}{@{}l}{\emph{Panel B. Full orders held fixed (Proposition~\ref{prop:control})}}\\
Entry \(\Ent\) & \reg{p4_frozen_r0_entry} & \reg{p4_frozen_r1_entry} & \\
High-class ownership \(\Own_H\) & \reg{p4_frozen_r0_own} & \reg{p4_frozen_r1_own} & \\
Expected target proceeds & \reg{p4_frozen_r0_proceeds} & \reg{p4_frozen_r1_proceeds} & \\
\midrule
\multicolumn{4}{@{}l}{\emph{Panel C. The challenger does not see the price}}\\
Orders \((\sigma_H,\sigma_L)\) & \((\delta_0,\delta_0)\) & \((\delta_1,\delta_{-1})\) & \\
Entry \(\Ent\) & \reg{p4_hidden_r0_entry} & \reg{p4_hidden_r1_entry} & \\
High-class ownership \(\Own_H\) & \reg{p4_hidden_r0_own} & \reg{p4_hidden_r1_own} & \\
Expected target proceeds & \reg{p4_hidden_r0_proceeds} & \reg{p4_hidden_r1_proceeds} & \\
\bottomrule
\end{tabular}
\par\medskip
\begin{minipage}{0.92\linewidth}\footnotesize
\textit{Notes.} Benchmark \(B\). Panel A: the unique equilibrium at each strength. Panel B: full orders imposed at
both strengths; at \(r_0\) this profile is a control, not an equilibrium. Panel C: the unique equilibrium when the
challenger decides on its prior. At \(r_2\) the full-order profile is the equilibrium of Panel A.
\end{minipage}
\end{table}

Table~\ref{tab:outcomes} reports both controls. Panel~B holds the full orders fixed. Only the top posterior
can recruit, because the threshold \(\tau(r)\) exceeds the prior at every strength. Entry is therefore
\(\reg{p4_frozen_r0_entry}\) at every strength up to the ceiling \(r_C=\reg{rC_B}\), and \(\rho\) above it. From
\(r_0\) to \(r_1\) the probability in Proposition~\ref{prop:control} is \(\reg{p4_control_crossing_r0_r1}\); from
\(r_1\) to \(r_2\) it is \(\reg{p4_control_crossing_r1_r2}\). With the information held fixed, entry is unchanged
from \(r_0\) to \(r_1\), and it falls once the ceiling binds. Full orders are not an equilibrium at \(r_0\), since a
purchase earns less than its cost there; this row is a control.

Panel~C hides the price from the challenger. The challenger decides on its prior and prepares only at zero cost.
Entry therefore stays at \(\reg{p4_hidden_r0_entry}\) at both strengths, although the investor trades fully at \(r_1\).
Together, the two controls locate the reversal. Incumbent strength changes the information that trading
generates, and the entry response to that information exceeds the direct deterrence effect.

\subsection{Coexistence}\label{sec:coexistence}

Proposition~\ref{prop:ccc} compares \(r_0\) with \(r_1\). Between them, informative and uninformative equilibria
coexist. This section gives the complete equilibrium set along the benchmark path \(r\in[r_0,r_2]\) at the prior
\(1/2\), for every slot in the band of Proposition~\ref{prop:region}(a). Write \(M=(1+u)/(2+u)\) and \(m=1-M\).

Two functions organize the path. The first is \(J(r)=\DT(r)[1-\tau(r)]\), the gross return of a purchase that just
reaches the threshold and recruits. The second is \(\fr(d)=\ell+d+\sqrt{d^2+2\ell d}\), the root above \(\ell\) of
\(\DT(r)=d\). Six strengths mark the changes in the equilibrium set; \(r_A\), \(r_N\) and \(r_J\) do not depend on
the slot.
\begin{center}\small\singlespacing
\begin{tabular}{@{}lll@{}}
\toprule
Strength & Definition & Role\\
\midrule
\(r_A\) & \(J(r)=k\) & a purchase that just reaches \(\tau(r)\) and recruits earns \(k\)\\
\(r_F\) & \(\fr\bigl(k/(1-M)\bigr)\) & a full purchase that recruits earns \(\DT(1-M)=k\)\\
\(r_N\) & \(\fr(2k/\rho)\) & against no informed trade, a unit order earns \(\rho\DT/2=k\)\\
\(r_J\) & \(\rho J(r)=k\) & a purchase that just reaches \(\tau(r)\) without recruitment earns \(k\)\\
\(r_S\) & \(\fr\bigl(k/(\rho m)\bigr)\) & a full sale earns \(\rho\DT m=k\)\\
\(r_C\) & \(\tau(r)=M\) & above it, no order recruits\\
\bottomrule
\end{tabular}
\end{center}

Write \(s_+=\lambda\sigma_H(1)\) for the probability of an informed purchase and \(s_-=\lambda\sigma_L(-1)\) for
that of an informed sale. No equilibrium uses a wrong-signed order (Lemma~\ref{lem:trading}(a)). Five forms of the
high-class order appear on the path:
\begin{center}\small\singlespacing
\begin{tabular}{@{}lll@{}}
\toprule
Form & Purchase probability \(s_+\) & Entry after a purchase\\
\midrule
N, no informed trade & \(0\) & \(\rho\)\\
R, a mixed purchase that recruits & \(b\,[\DT(r)/k-2]\) & one\\
F, a full purchase & \(\lambda\) & one up to \(r_C\), \(\rho\) above\\
Q, mixed orders at the entry floor & \(b\,[\rho\DT(r)/k-2]\) & \(\rho\)\\
T, the threshold pair & \(b\,(2\tau-1)/(1-\tau)\), so \(\mu(1)=\tau\) & \(k/J(r)\)\\
\bottomrule
\end{tabular}
\end{center}
In each form the low class sells with probability \(s_-=0\) on \([r_0,r_N]\), \(s_-=b\,[\rho\DT(r)/k-2]\) on
\((r_N,r_S)\), and \(s_-=\lambda\) on \([r_S,r_2]\). Entry after no order and after a sale is \(\rho\), and the
prices follow from equation~\eqref{eq:price}. Each form thus defines a complete tuple.

\begin{proposition}[The equilibrium set on the benchmark path; computer-assisted]\label{prop:coexistence}
Take the benchmark auction, the prior \(1/2\) and any slot in the band of Proposition~\ref{prop:region}(a). Then
\(r_0<r_A<r_F<r_N<r_J<r_S<r_1<r_C<r_2\), and the equilibrium set at each strength is
\begin{center}
\begin{tabular}{@{}ll@{}}
\toprule
Strength & Equilibria\\
\midrule
\(r_0\le r<r_A\) & N\\
\(r=r_A\) & N; T with entry one after a purchase (the left end of R)\\
\(r_A<r<r_F\) & N; R\\
\(r_F\le r\le r_N\) & N; F\\
\(r_N<r<r_J\) & F; Q\\
\(r_J\le r\le r_2\) & F\\
\bottomrule
\end{tabular}
\end{center}
Entry is \(\rho\) in N and Q, \(\Ent_R(r)=\rho+(1-\rho)\,b\DT(r)/(2k)\) in R, and
\(\Ent_F=\rho+(1-\rho)(b+\lambda/2)\) in F up to \(r_C\), with \(\Ent_F=\rho\) above \(r_C\).
\begin{enumerate}[label=(\alph*),leftmargin=2em]
\item No informed trade and an informative equilibrium coexist exactly on \([r_A,r_N]\).
\item On the informative branch that coexists with no trade, entry is
\(\rho+(1-\rho)\min\{b\DT(r)/(2k),\,b+\lambda/2\}\). It is continuous, rises strictly on \([r_A,r_F]\), and is
constant on \([r_F,r_N]\). It keeps this value along F up to \(r_C\) and equals \(\rho\) above \(r_C\).
\item The equilibrium set is nonempty at every strength.
\end{enumerate}
\end{proposition}

The status is computer-assisted because two monotonicity steps rest on exact Sturm counts: \(J\) and
\(\DT(1-\tau)^2\) rise along the path. Every other step is analytical. At the benchmark slot,
\(r_A=\reg{rA}\), \(r_F=\reg{rF_B}\), \(r_N=\reg{rN}\), \(r_J=\reg{rJ}\), \(r_S=\reg{rS_B}\) and
\(r_C=\reg{rC_B}\). On R, entry is \(\reg{entry_rA_B}\) at \(r_A\), \(\reg{entry_mb_r140_B}\) at
\(\reg{str_r140}\), and \(\reg{entry_mb_r145_B}\) at \(\reg{str_r145}\). From \(r_F\) it stays at
\(\reg{entry_r1_B}\) until \(r_C\). The low class abstains up to \(r_N\), mixes on \((r_N,r_S)\), and sells with
probability one from \(r_S\).

\begin{figure}[tbp]
\centering
\slotfigure{figures/coexistence.pdf}{7cm}
\caption{Equilibrium entry along the benchmark path.}\label{fig:coexistence}
\par\medskip
\begin{minipage}{\linewidth}\footnotesize\singlespacing
\textit{Notes.} The benchmark \(B\) at the prior \(1/2\). Entry against incumbent strength \(r\) in each equilibrium
form of Proposition~\ref{prop:coexistence} under rules IM and C1. The thin dashed line is the threshold pair T, an
equilibrium under rule CM only. Panel (a) covers the path from \(r_0\) to \(r_2\), and panel (b) enlarges the
strengths around the coexistence interval. On F the high class buys fully; the low class abstains up to \(r_N\),
mixes from \(r_N\) to \(r_S\), and sells fully from \(r_S\). The blue band from \(r_A\) to \(r_N\) marks
coexistence. Filled markers belong to a branch, and open markers do not. Above \(r_C\), entry on F equals \(\rho\).
\end{minipage}
\end{figure}

Figure~\ref{fig:coexistence} shows entry along the path. Against an uninformative price with entry \(\rho\), a
unit order earns \(\rho\DT/2\), which stays below \(k\) up to \(r_N\), so no trade survives until \(r_N\). An
informative equilibrium changes the entry response and with it the return to trade, so it can hold at the same
strengths.

Entry rises strictly along R because the purchase rate adjusts. At \(r_A\) a purchase that just reaches the
threshold earns exactly \(k\). Above \(r_A\), class \(H\) mixes, and the purchase posterior stays above the
threshold. The purchase then earns \(\DT[1-\mu(1)]=k\), so a larger spread needs a more informative purchase. The
purchase rate \(s_+=b\,[\DT(r)/k-2]\) therefore rises with \(r\), and entry rises with it. At \(r_F\) the purchase
rate reaches \(\lambda\), and from there entry is flat. Entry therefore rises weakly along the informative branch,
and strictly where the investor mixes.

Above \(r_N\), no trade fails. On \((r_N,r_J)\) the investor can also mix in both directions without recruitment
(Q), and from \(r_J\) on, full purchases are the only high-class choice.

\begin{remark}[Other tie rules]\label{rem:ties}
Under the challenger rule CM, the threshold pair T is also an equilibrium on \((r_A,r_J]\). Its entry falls strictly
from \(\Ent_R(r_A)\) at \(r_A\) to \(\rho\) at \(r_J\); at the benchmark it is \(\reg{entry_th_r155_B}\),
\(\reg{entry_th_r160_B}\) and \(\reg{entry_th_r165_B}\) at \(\reg{str_r155}\), \(\reg{str_r160}\) and
\(\reg{str_r165}\). At \(r_C\), entry after a purchase can then take any value in \([\rho,1]\). 

Under the investor rule I0, the equilibrium set is N on \([r_0,r_F]\); N and F on \((r_F,r_N]\); empty on \((r_N,r_S]\); and F on
\((r_S,r_2]\). On \((r_N,r_S]\) the low class has no admissible sale rate under I0. Abstention fails, because a
sale earns more than \(k\). A full sale fails, because it earns at most \(k\). Under the investor rule IM, the
equilibrium set is nonempty along the whole path. The strict rise of entry along R also needs IM.
\end{remark}

Coexistence also holds at other priors. For a prior \(\pi\), write \(r_N(\pi)=\fr\bigl(k/(\rho\max\{\pi,1-\pi\})\bigr)\).
When some strength on the path has \(\tau(r)\le M(\pi)\), let \(r_C(\pi)\) be the largest one. When
\(\tau(r_0)\le\pi\le\tau(r_2)\), let \(r_\pi\) solve \(\tau(r_\pi)=\pi\). Three priors divide the cases: \(\pi_A^{\mathrm{lo}}\) solves
\(M(\pi)=\tau(r_A)\); \(\pi_A^{\mathrm{hi}}=\tau(r_A)\); and \(\pi_{\max}=\tau(r_*)\), where \(r_*\) solves
\(\rho\DT(r)\tau(r)=k\).

\begin{corollary}[Coexistence at every prior; computer-assisted]\label{prop:priors}
Take the benchmark auction, any slot, and any prior \(\pi\in(0,1)\). The set \(C(\pi)\) of strengths on the path
at which no informed trade and an informative equilibrium coexist is
\[
C(\pi)=\bigl\{r\in[r_0,r_2]:\ r\ge r_A,\ \ \pi<\tau(r)\le M(\pi),\ \ \rho\DT(r)\max\{\pi,1-\pi\}\le k\bigr\}.
\]
\begin{enumerate}[label=(\alph*),leftmargin=2em]
\item If \(\pi<\pi_A^{\mathrm{lo}}\), \(C(\pi)\) is empty. At \(\pi=\pi_A^{\mathrm{lo}}\), \(C(\pi)=\{r_A\}\).
\item If \(\pi_A^{\mathrm{lo}}<\pi<\pi_A^{\mathrm{hi}}\), then \(C(\pi)=[r_A,\min\{r_N(\pi),r_C(\pi)\}]\), a
nondegenerate interval. Coexistence starts at \(r_A\) for every prior in this window and every slot.
\item If \(\pi_A^{\mathrm{hi}}\le\pi<\pi_{\max}\), then \(C(\pi)=(r_\pi,\min\{r_N(\pi),r_C(\pi)\}]\), a nonempty
interval with \(r_\pi\ge r_A\). Coexistence starts where the prior reaches the threshold.
\item If \(\pi\ge\pi_{\max}\), \(C(\pi)\) is empty.
\end{enumerate}
At each \(r\in C(\pi)\), the informative equilibrium is R or F (the threshold pair at \(r_A\) in case (b)). Its
entry is \(\rho+(1-\rho)\min\{(1-\pi)\,b\DT(r)/k,\ b+\pi\lambda\}\), which rises strictly while the purchase is
mixed and is then constant.
\end{corollary}

At the benchmark slot, \(\pi_A^{\mathrm{lo}}\approx\reg{piA_lo_B}\), \(\pi_A^{\mathrm{hi}}\approx\reg{piA_hi}\) and
\(\pi_{\max}\approx\reg{pimax}\). For priors between \(\pi_A^{\mathrm{lo}}\) and \(\pi_A^{\mathrm{hi}}\), the
informative branch starts at the same strength \(r_A\), where both the investor and the challenger are indifferent.
Between \(\pi_A^{\mathrm{hi}}\) and \(\pi_{\max}\), it starts later, at \(r_\pi\). Below \(\pi_A^{\mathrm{lo}}\), even the
top posterior cannot reach the threshold once a purchase can pay for itself. From \(\pi_{\max}\) on, the
uninformative equilibrium fails before the prior falls below the threshold. Under rule CM, a challenger tie at the prior adds one strength: for every
\(\pi\in(\pi_A^{\mathrm{hi}},\pi_{\max}]\), the set \(C(\pi)\) also contains \(r_\pi\). At \(\pi=\pi_{\max}\) this
point is the whole coexistence set.
